\documentclass{jsarticle}
\usepackage{amsmath}
\usepackage{amssymb}
\begin{document}
\title{不完全性定理と完全数、\\
そして自然法則とそれを表す数式群}
\author{山口 雅旭   緒方一之}
\date{}
\maketitle

  不完全性定理は、虚数と平方根が$ e^{ix},\sqrt{x} $により、
  $ 2^0,1^0,x^0 = 1 $と合わせて,記号とそれを説明している自然の法則には、
  ギャップが存在していて、このギャップが数式の記号を仮定として、
  数学の決まり事として自然と数学の近似式として、
  完全な自然を説明する道具が無く、そのために完全な神は存在していないと
  ゲーデルは証明している。つまり、数学の決まり事と自然の法則に
  照らし合わされている。そのために、記号と自然法則に隔たりが生じている。
  自然法則に数式は近似と仮定として成り立っている。そのために、
  絶対神は存在していないとゲーテルは証明した。 
 
  しかし、ゼータ関数により、その数式を解析するのために
  使うマインドマップによれば、偏微分と多重積分の差集合から、
  非可換代数方程式に求めれば、
  $$ \pi(\chi,x)  = [i \pi(\chi,x),f(x)] $$
  これは世界線として、２次曲面を共変微分を大域的な領域で接平面で切って、
  質量がヒッグス場で時間と空間の関係として表されている。
  $$ {d \over df}F = {d \over df}\int \int{1 \over (x \log x)^2}dx_m
  + {d \over df}\int \int{1 \over (y \log y)^{1 \over 2}}dy_m $$
  $$ ||ds^2|| = e^{-2 \pi T |\phi|}[\eta + \bar{h}_{\mu\nu}]dx^{\mu}
  dx^{\nu} + T^2 d^2 \psi $$

  この神が居ないとする不完全性定理を覆すのに、
  積分と数列の和の極限の機能の結果、このギャップを補助するこの差分方程式
  から、商代数から単体量が不変式としての結果であり、
  そして、不完全性定理がゼータ関数と量子群により、
  記号と自然法則の完全数が求まり、原子と電子、そして素粒子がこのラプラス
  方程式から宇宙のデータベースにアクセスして、ホログラム空間として、
  記号と自然法則の写像として、不完全性定理の意味と法則が書き換えられる。


  ブルバギでは、写像と関数を使うと、近似式は関数で、写像を使うと近似式を
  完全数として書き換えられる。不完全性定理はトポロジーを使うと記号と数式の
  自然法則を完全数として破られる。







\end{document}
